% =====================================================================
%  Capítulo 1. Introducción
% =====================================================================
\chapter[INTRODUCCIÓN]{Introducción}
En las últimas décadas, las Interfaces Cerebro Computadora (BCI, por sus siglas en inglés) se
han convertido en una de las áreas de investigación más relevantes dentro de la biónica y la
inteligencia artificial orientada al área de la salud, al permitir entablar comunicación directa
entre la actividad cerebral de un individuo y un dispositivo externo sin depender de ninguna
señal o movimiento muscular. Como se menciona en el trabajo de Felton et al. \cite{felton2007},
las BCI traducen señales cerebrales registradas en acciones en tiempo real, lo cual abre la
posibilidad de restaurar funciones comunicativas o motoras en personas con discapacidad severa.

Dentro de los paradigmas más estudiados en BCI basadas en electroencefalografía (EEG) se
encuentra el uso de los Potenciales Relacionados con Eventos (ERP), específicamente la
componente P300, cuya detección mediante la iluminación intermitente de filas y columnas de
una matriz de caracteres dio origen al llamado P300 Speller \cite{farwell1988}. A partir de este
paradigma se han desarrollado sistemas de comunicación asistida tanto en el ámbito nacional
como internacional, alcanzando en algunos casos porcentajes de acierto de hasta 95\% en
evaluaciones intrasujeto y tasas sostenidas superiores a 17 caracteres por minuto cuando se
emplean señales invasivas de electrocorticografía (ECoG). Sin embargo, de estos antecedentes se
identifican dos limitaciones persistentes en los sistemas P300 basados en EEG: la necesidad de
múltiples canales para mantener niveles altos de precisión, y la dependencia de múltiples
repeticiones de estímulos para asegurar la detección confiable del componente P300, lo que
reduce la velocidad de comunicación y aumenta la carga cognitiva del usuario.

Bajo este contexto, el presente Trabajo Terminal se orienta al desarrollo de un sistema de
comunicación asistida basado en el paradigma P300 Speller utilizando EEG, integrando modelos
de aprendizaje profundo para la detección automática del componente P300 e incorporando
estrategias para reducir el número de canales y el número de repeticiones requeridas por
selección de carácter, con el objetivo de incrementar la velocidad de comunicación sin
comprometer la precisión del sistema.

\section[OBJETIVOS]{Objetivos}

\subsection{Objetivo general}

Desarrollar un sistema de comunicación asistida implementado en Python que permita a las
personas con problemas motrices severos comunicarse por medio de señales
electroencefalográficas usando el paradigma P300 Speller, a través de una interfaz gráfica
basada en una matriz de selección visual de caracteres mediante la iluminación intermitente de
filas y columnas, clasificando las componentes P300 de las señales cerebrales utilizando un
modelo de aprendizaje profundo desarrollado en Python.

\subsection{Objetivos específicos}

\begin{enumerate}
    \item Analizar y recolectar bases de datos públicas de señales electroencefalográficas
    sobre el paradigma P300 Speller, con el fin de identificar las características más
    importantes para el desempeño del modelo.
    \item Revisar el estado del arte para seleccionar y analizar los modelos de aprendizaje
    profundo más adecuados para la detección de la componente P300 en señales cerebrales.
    \item Diseñar e implementar un modelo de aprendizaje profundo para la clasificación de la
    componente P300 en señales EEG.
    \item Evaluar el desempeño del modelo mediante métricas de clasificación buscando minimizar
    los canales o ciclos necesarios para la detección de la componente P300.
    \item Diseñar e implementar el sistema de comunicación asistida, integrando el modelo de
    detección de P300 con una interfaz que permita la construcción de palabras por parte del
    usuario.
    \item Integrar un módulo de predicción de palabras o autocompletado, orientado a reducir el
    número de selecciones necesarias para la construcción de mensajes.
    \item Evaluar el funcionamiento del sistema mediante pruebas que permitan analizar su
    desempeño en la selección de caracteres y construcción de palabras.
\end{enumerate}

\section[JUSTIFICACIÓN]{Justificación}

El análisis de señales EEG en paradigmas basados en la componente P300 ha sido ampliamente
utilizado en el desarrollo de interfaces cerebro-computadora, particularmente en sistemas de
selección de caracteres mediante el paradigma P300 Speller. Estos sistemas permiten la
comunicación mediante la detección de respuestas cerebrales asociadas a estímulos visuales, lo
que representa una alternativa para las personas con compromisos motores severos que tengan
problemas de comunicación.

De acuerdo con el trabajo de Feigin et al. \cite{feigin2020}, los desórdenes neurológicos son la
principal causa de discapacidad a nivel mundial, y es posible que la magnitud de estas aumente
debido al crecimiento y envejecimiento de la población; en enfermedades como la esclerosis
lateral amiotrófica (ELA) la comunicación puede verse gravemente limitada, lo que impacta
directamente en la calidad de vida del paciente.

Dentro de este contexto, el paradigma P300 Speller puede resultar útil para permitir la
comunicación a las personas con problemas motrices severos, mejorando así su calidad de vida.
Sin embargo, la mayoría de los estudios se enfocan en maximizar la precisión de clasificación
mediante entornos controlados difíciles de llevar a la práctica, de modo que este trabajo
interviene en la integración estructural del sistema, priorizando no solo el desempeño del
modelo, sino también su incorporación en un sistema funcional que le permita a los usuarios
objetivo la construcción efectiva de mensajes.

La reducción del número de canales o ensayos, así como la incorporación de estrategias que
faciliten la selección de caracteres, puede contribuir al desarrollo de sistemas BCI más
accesibles, portables y de implementación simplificada, ampliando su viabilidad en distintos
contextos tecnológicos y acercando su uso a escenarios reales de asistencia comunicativa. De
esta manera, el proyecto no se centra solamente en la optimización bajo escenarios aislados,
sino que propone un sistema orientado a la aplicación práctica.

En cuanto a la aportación del trabajo, su valor radica en la integración de un modelo de
aprendizaje profundo para la detección de la componente P300 con un modelo de predicción de
caracteres, combinación que permite no solo identificar la selección de caracteres mediante las
señales obtenidas por electroencefalografía, sino también optimizar el proceso de construcción
de mensajes mediante la asistencia de este otro módulo.

Por lo tanto, la complejidad del trabajo implica el procesamiento y análisis de señales EEG, la
implementación y entrenamiento de modelos de aprendizaje profundo, la integración de estos
modelos en una aplicación y la evaluación de los resultados, lo cual requiere la aplicación
avanzada de conocimientos en procesamiento de señales, aprendizaje profundo y desarrollo de
sistemas, competencias propias del área de ingeniería en inteligencia artificial. Finalmente, el
proyecto resulta viable en términos de tiempo y recursos, ya que los modelos de aprendizaje
profundo pueden entrenarse utilizando bases de datos públicas de paradigmas P300, evitando la
necesidad de adquisición propia de las mismas y reduciendo costos asociados a hardware
especializado con 64 canales de detección de señales EEG.
